\scanpage{50}
% archive sheet number: 49
groupe des autom. de \((X,\Delta)\) qui fixe \(\infty\) et \(-1\) donc \inserted{des \unsure{autos|autom.}} de la forme \(z\mapsto az+b\) avec \(-a+b=-1\)
\[
\text{i.e.}\qquad z\longmapsto\deleted{a(z-1)+1=az+(1-a)}az+(a-1)
\]
est plus petit, en fait si l'autom envisagé n'est pas l'identité, il doit induire une permutation sans pts fixes de \((0,1,i,-i)\), qui est donc d'ordre \(2\) ou quatre, donc on a \(a^4=1\) \deleted{ou \(a^4=1\)} \deleted{i.e. \(a\in\{0,1,i,\unsure{-}\)} mais faisant \(z=0\) on doit avoir
\[
a-1\in\{1,i,-i\}\qquad\text{i.e.}\quad a\in(2,1+i,1-i)
\]
% slanted
\margin{donc \(a^4=\pm4\) qui ne peut être \(=1\) qu'en car \(3\) et \(5\)}
faisant \(z=1\) on trouve
\[
2a-1\in\{0,i,-i\}\qquad\text{i.e.}\quad
a\in\{1/2,\tfrac12(i+1),\tfrac12(1-i)\}
\]
% slanted
\margin{donc \(a^4=1/4\) ou \(a^4=-1\) qui n'est pas égal à \(1\) (car on est en car \(\ne2\))}
\(\bigl[\)faisant \(z=i\), on trouve
\[
a\underbrace{(i+1)}_{1:\left(\frac{1-i}{2}\right)}-1\in\{0,1,-i\}\qquad\text{i.e.}\quad
a\in\left\{\frac{1\deleted{\unsure{+}}\unsure{-}i}{\deleted{\unsure{4}}2},\deleted{\unsure{\frac{2}{1+i}}}1-i,-i\right\}
\]
et de \(=\), faisant \(z=-i\), on trouve
\[
a\in\left\{\frac{1+i}{2},1+i,i\right\}\biggr]
\]
Donc l'\deleted{\unsure{a}}intersection des stabilisateurs de \(\infty,-1\) est \(\{1\}\), et par raison d'homogénéité il en est de \(=\) pour \deleted{\unsure{tout|toute}} partie à deux éléments de \(\Delta\), distincte de \(\Delta_1,\Delta_2,\Delta_3\). Donc on voit que la structure octaédrale \(\Delta=\Delta_1\amalg\Delta_2\amalg\Delta_3\) est invariante par \(\Gamma\), et uniquement déterminée par le diviseur \(\Delta\). Donc \(\Gamma\hookrightarrow\underline{\operatorname{Aut}}_{\mathrm{oct}}(\Delta)\),
% slanted
\margin{et on a \(=\)\\ \(\Gamma\subset\operatorname{Aut}^{+}_{\mathrm{oct}}(\Delta)\)}
Pour montrer qu'on a \unsure{isomorphisme|isomorphie}, il suffit de noter que le groupe des autom \inserted{directs} de la structure octaédrale est engendrée par les trois sous-groupes cycliques d'ordre \(4\) \deleted{\unsure{\(t_1=\)}} qu'il contient, correspondants à \(\Delta_1,\Delta_2,\Delta_3\) -- or ceux-ci sont contenus dans \(\Gamma\).
